% Part: normal-modal-logic
% Chapter: tableaux
% Section: introduction

\documentclass[../../../include/open-logic-section]{subfiles}

\begin{document}

\olfileid{nml}{tab}{int}

\olsection{ভূমিকা}

!!^{tableau}s হলো !!{signed
  formula}s-এর এক ধরনের (নিচের দিকে
শাখাবিশিষ্ট) বৃক্ষ। প্রতিটি চিহ্নিত সূত্রে একটি সত্যমানের চিহ্ন
($\True$ বা $\False$) এবং !!a{sentence} থাকে:
\[
\sFmla{\True}{!A} \text{ অথবা } \sFmla{\False}{!A}.
\]
!!^a{tableau} কয়েকটি \emph{অনুমিতি} দিয়ে শুরু হয়। পরবর্তী প্রতিটি
!!{signed formula} কোনো অনুমানবিধি প্রয়োগ করে পাওয়া যায়। কিছু বিধি
বৃক্ষের একটি অগ্রভাগে এক বা একাধিক !!{signed formula}s যোগ করে; অন্যগুলি
দুটি নতুন অগ্রভাগ সৃষ্টি করে, ফলে দুটি শাখা হয়। যে !!{signed formula}-তে
বিধি প্রয়োগ করা হয়, বিধি থেকে পাওয়া !!{signed formula}s-এর মূল !!{formula}
তার চেয়ে কম জটিল। কোনো শাখায় $\sFmla{\True}{!A}$ ও
$\sFmla{\False}{!A}$ উভয়ই থাকলে শাখাটিকে \emph{বদ্ধ} বলি। !!a{tableau}-র
প্রতিটি শাখা বদ্ধ হলে গোটা !!{tableau} বদ্ধ। একটি বদ্ধ !!{tableau}
হলো এমন !!a{derivation}, যা দেখায় যে !!{tableau} শুরু করার জন্য ব্যবহৃত
!!{signed formula}s-এর সমষ্টি অপরিতৃপ্তিযোগ্য। এর সাহায্যে $\Proves$ সম্পর্ক
সংজ্ঞায়িত করা যায়: $\Gamma \Proves !A$ তখনই এবং কেবল তখনই, যখন
কোনো সসীম সমষ্টি~$\Gamma_0 = \{!B_1,
\dots, !B_n\} \subseteq \Gamma$-র জন্য নিচের অনুমিতিগুলির একটি বদ্ধ !!{tableau} আছে:
\[
\{\sFmla{\False}{!A}, \sFmla{\True}{!B_1}, \dots, \sFmla{\True}{!B_n}\}.
\]

মোডাল যুক্তিবিদ্যার ক্ষেত্রে !!{signed
formula}-র ধারণা প্রসারিত
করতে এবং~\iftag{prvBox}{$\Box$\iftag{prvDiamond}{ ও
    $\Diamond$}}{$\Diamond$}-এর জন্য বিধি যোগ করতে হয়। মোডাল
!!{tableau}s-এর !!{formula}s-এ চিহ্নের ($\True$ বা $\False$) পাশাপাশি
\emph{পূর্বসূচি}~$\sigma$ থাকে। পূর্বসূচি হলো ধনাত্মক পূর্ণসংখ্যার
অশূন্য অনুক্রম, অর্থাৎ $\sigma \in (\PosInt)^* \setminus
\{\emptyseq\}$। লেখার সময় আমরা পূর্বসূচির চারপাশে
$\tuple{\ }$ রাখি না এবং !!{element}s-কে $,$-র বদলে $.$ দিয়ে
আলাদা করি। $\sigma$ পূর্বসূচি হলে $\sigma.n$ হলো $\sigma
\concat \tuple{n}$; যেমন, $\sigma = 1.2.1$ হলে $\sigma.3$ হলো
$1.2.1.3$। অতএব, উদাহরণস্বরূপ,
\[
\sFmla{\True}{\Box !A \lif !A}[1.2]
\]
একটি \emph{পূর্বসূচিযুক্ত !!{signed formula}} (সংক্ষেপে,
\emph{পূর্বসূচিযুক্ত !!{formula}})।

স্বতঃবোধ্য অর্থে, পূর্বসূচিটি মডেলের এমন একটি জগতের নাম দেয় যেখানে
!!a{tableau}-র কোনো শাখার !!{formula}s সত্য হতে পারে। $\sigma$ যদি
কোনো জগতের নাম দেয়, তবে $\sigma.n$, $\sigma$-র নামাঙ্কিত জগৎ থেকে অভিগম্য
একটি জগতের নাম দেয়।

\end{document}
